\documentclass[11pt]{article}
\usepackage[utf8]{inputenc}
\usepackage[english]{babel}
\usepackage{geometry}
\usepackage{amsmath}
\usepackage{booktabs}
\usepackage{xcolor}
\usepackage{mdframed}

\geometry{a4paper, margin=1in}

\title{Technical Specifications: HackRF One}
\author{Radiofrequency Engineering}
\date{\today}

\newmdenv[
  backgroundcolor=red!8,
  linecolor=red!60,
  linewidth=1pt,
  roundcorner=4pt,
  innertopmargin=6pt,
  innerbottommargin=6pt
]{warningbox}

\newmdenv[
  backgroundcolor=yellow!10,
  linecolor=yellow!60!black,
  linewidth=1pt,
  roundcorner=4pt,
  innertopmargin=6pt,
  innerbottommargin=6pt
]{pendingbox}

\begin{document}

\maketitle
\tableofcontents
\newpage

% =========================================================
\section{Gains}
% =========================================================

Controlled separately through \texttt{libhackrf} (tools such as \texttt{hackrf\_transfer} or GNU Radio):

\begin{itemize}
    \item \textbf{LNA (IF RX):} \texttt{hackrf\_set\_lna\_gain()} - 0 to 40 dB, adjustable in 8 dB steps. Affects the signal in the intermediate-frequency (IF) stage.
    \item \textbf{VGA (Baseband RX):} \texttt{hackrf\_set\_vga\_gain()} - 0 to 62 dB, adjustable in 2 dB steps. Variable-gain amplifier in baseband.
    \item \textbf{TX VGA:} \texttt{hackrf\_set\_txvga\_gain()} - 0 to 47 dB, adjustable in 1 dB steps.
    \item \textbf{AMP (RF Front End):} \texttt{hackrf\_set\_amp\_enable()} - 0 (off) or 1 (on). Enables the RF front-end amplifier, providing a fixed $\approx$+14 dB.
    \begin{pendingbox}
    The amplifier part number varies depending on the hardware revision. The MGA-81563 appears in third-party schematics for early revisions, but it is not confirmed as an official specification in the Great Scott Gadgets documentation.
    \end{pendingbox}
    Applies to the final TX stage or the initial RX stage.
\end{itemize}

\textbf{Note:} Since the ADC is 8-bit (dynamic range limited to $\approx$48 dB), correct gain management is crucial to avoid raising the noise floor or causing clipping during analog-to-digital conversion.

% =========================================================
\section{Hardware Filters and Processing}
% =========================================================

\begin{itemize}
    \item \textbf{Analog Filter (Baseband LPF):} \texttt{hackrf\_set\_baseband\_filter\_bandwidth()}\\
    Selectable thresholds in the MAX2837 (r1--r8) / MAX2839 (r9) transceiver: 1.75, 2.5, 3.5, 5, 5.5, 14, 20, 24, and 28 MHz. Helps prevent analog aliasing.
    \item \textbf{CPLD and Microcontroller (No Internal DSP):}\\
    Uses a Xilinx XC2C64A CPLD for signal routing and an ARM Cortex-M4 MCU (NXP LPC4320/30). Unlike other high-end SDRs, the HackRF One \textbf{does not perform internal digital signal processing} (such as DDC/DUC or hardware decimation). All mathematical processing falls to the host CPU over USB.
\end{itemize}

% =========================================================
\section{Frequency Ranges}
% =========================================================

\begin{itemize}
    \item \textbf{Nominal Range:} 1 MHz to 6 GHz.
    \item \textbf{Synthesis Architecture:}
    \begin{itemize}
        \item 2.3 GHz to 2.7 GHz: Direct conversion using the MAX2837/MAX2839 transceiver.
        \item Outside that band: Uses the RFFC5071 mixer to perform up-conversion or down-conversion into the transceiver band.
    \end{itemize}
    \item \textbf{Performance note:} Below 10 MHz, the efficiency of the RF transformers (baluns) decreases.
\end{itemize}

% =========================================================
\section{Sample Rate and Resolution}
% =========================================================

Limited by the ADC/DAC (MAX5864) capability and the USB 2.0 bus.

\begin{table}[h]
    \centering
    \begin{tabular}{llll}
        \toprule
        \textbf{Mode} & \textbf{I/Q Resolution} & \textbf{Sample Rate} & \textbf{Maximum Bandwidth} \\
        \midrule
        Standard & 8-bit I / 8-bit Q & 2 to 20 MS/s & 20 MHz (theoretical Nyquist) \\
        \bottomrule
    \end{tabular}
\end{table}

\textbf{Note:} To avoid packet drops (dropped samples) at 20 MS/s, it is recommended to connect the HackRF to a dedicated USB controller on the host that does not share bandwidth with other high-demand peripherals. Rates below 2 MS/s are not officially supported.

% =========================================================
\section{IQ \& DC Corrections}
% =========================================================

\begin{itemize}
    \item \textbf{DC Peak (DC Offset / LO Leakage):} Due to its homodyne architecture (Zero-IF / Direct Conversion), the HackRF One shows a visible spike at the center of the spectrum (0 Hz in baseband). It \textbf{must be corrected in software} (e.g., \textit{DC Blocker} in GNU Radio) or by using \textit{Offset Tuning} (tuning a few MHz to the side and digitally shifting).
    \item \textbf{I/Q Imbalance:}
    \begin{pendingbox}
    It can appear as ghost mirror frequencies when working with strong signals. It is recommended to calibrate IQ gain and phase in the host software (SDR\#, GNU Radio), since the device does not apply automatic factory corrections.
    \end{pendingbox}
\end{itemize}

% =========================================================
\section{Connectors and Inputs}
% =========================================================

\begin{table}[h]
    \centering
    \small
    \begin{tabular}{p{0.18\textwidth}p{0.26\textwidth}p{0.44\textwidth}}
        \toprule
        \textbf{Connector} & \textbf{Function} & \textbf{Specification} \\
        \midrule
        Micro-USB   & Power and data        & USB 2.0 High Speed (480 Mbps). 5 V / $\approx$500 mA. \\
        ANT (SMA)   & Main TX/RX            & Female, 50 $\Omega$. Half-duplex. 1 MHz -- 6 GHz. \\
        CLKIN (SMA) & External reference    & 10 MHz. Square wave, \textbf{0 V to 3.0 V}. Do not exceed 3.0 V or go below 0 V. Do not use other frequencies except with firmware modification. \\
        CLKOUT (SMA)& Output clock          & 10 MHz, 3.3 V. Enabled with \texttt{hackrf\_clock -o 1}. For daisy chaining between units. \\
        P20, P22, P28 & Expansion headers    & LPC4320 MCU pins. 3.3 V logic, for modules such as PortaPack. \\
        \bottomrule
    \end{tabular}
\end{table}

% =========================================================
\section{Safe Operating Limits}
% =========================================================

\begin{warningbox}
\textbf{Critical warnings - LNA / front-end susceptibility}
\end{warningbox}

\vspace{4pt}
\begin{itemize}
    \item \textbf{Maximum RX Power:} NEVER exceed \textbf{-5 dBm} ($\approx$0.3 mW) at the antenna port. Exceeding it will destroy the front-end amplifier. In theory, with the AMP disabled it may tolerate up to +10 dBm, but a software error that re-enables it would cause permanent damage - always use an external attenuator.
    \item \textbf{TX power by band} (reference for sizing loopback paths and amplifiers):
    \begin{itemize}
        \item 1--10 MHz: 5 to 15 dBm (increases with frequency).
        \item 10--2170 MHz: 5 to 15 dBm (decreases with frequency).
        \item 2170--2740 MHz: 13 to 15 dBm (best performance region).
        \item 2740--4000 MHz: 0 to 5 dBm.
        \item 4000--6000 MHz: $-$10 to 0 dBm.
    \end{itemize}
    \textbf{In TX$\to$RX loopback:} use sufficient attenuation for the band. Never connect TX directly to RX.
    \item \textbf{Unloaded RF port:} Do not transmit without an antenna or a 50 $\Omega$ dummy load. Reflected power can burn out the TX stage.
    \item \textbf{Clock Voltages:} Signals on CLKIN must not be negative or exceed \textbf{3.0 V}. A sine wave centered at 0 V requires DC offset to meet this requirement.
    \item \textbf{Antenna Power (Bias-T):} Can provide 3.3 V at 50 mA on the SMA connector to power external LNAs. Do not enable it if the connected device presents a DC short circuit.
\end{itemize}

% =========================================================
\section{Clock and Synchronization}
% =========================================================

\begin{itemize}
    \item \textbf{Internal Oscillator:} Standard 10 MHz crystal. Ppm stability is not officially specified by Great Scott Gadgets; community estimates suggest 20--30 ppm, susceptible to thermal drift.
    \item \textbf{TCXO Module (External/Additional):} For serious use at high frequencies or with narrow modulation, attach a 0.5 ppm TCXO module on header P22 or use CLKIN with a stable external reference.
    \item \textbf{CLKIN:}
    \begin{itemize}
        \item 10 MHz signal, 0 to 3.0 V square wave. High impedance.
        \item The device automatically switches to the external clock when the next TX/RX operation starts. Verify detection with \texttt{hackrf\_debug --si5351c -n 0 -r} (output \texttt{0x01} = clock detected).
        \item CLKOUT from one HackRF One can be directly daisy chained to CLKIN on another.
    \end{itemize}
    \item \textbf{CLKOUT:} Enable with \texttt{hackrf\_clock -o 1}. Allows synchronization of multi-HackRF One arrays with a single master clock.
\end{itemize}

% =========================================================
\section{Indicator LEDs}
% =========================================================

Six LEDs are mounted along the board. When connected to USB, \textbf{four should turn on}: 3V3, 1V8, RF, and USB.

\begin{table}[h]
    \centering
    \begin{tabular}{p{0.10\textwidth}p{0.38\textwidth}p{0.40\textwidth}}
        \toprule
        \textbf{LED} & \textbf{Function} & \textbf{State and interpretation} \\
        \midrule
        \textbf{3V3} & Main power & On when USB is connected. Indicates that the 3.3 V regulator is working. If it does not turn on: hardware fault. \\
        \addlinespace
        \textbf{1V8} & Active firmware + secondary power & On when USB is connected. Indicates that the firmware is running and has enabled additional internal power rails (1.8 V). If it does not turn on while 3V3 is on: firmware problem. \\
        \addlinespace
        \textbf{RF}  & RF chain powered & On when USB is connected. Indicates that the firmware has enabled power to the RF block. \\
        \addlinespace
        \textbf{USB} & Active USB communication & On when the host recognizes the device correctly. \\
        \addlinespace
        \textbf{RX}  & Reception in progress & On \textbf{only} during active captures. \\
        \addlinespace
        \textbf{TX}  & Transmission in progress & On \textbf{only} during active transmission. \\
        \bottomrule
    \end{tabular}
\end{table}

\textbf{Note:} Adjacent LED colors are different to make visual distinction easier, but the colors themselves have no functional meaning.

% =========================================================
\section{Expansion Pinout (P20 / P22 / P28)}
% =========================================================

Provides access to microcontroller GPIOs (designed to interconnect add-ons such as the PortaPack). They operate at \textbf{3.3 V}.

\subsection*{P20 and P28}
\begin{itemize}
    \item Contain pins for the LPC internal ADC, I2C, SPI, and auxiliary power.
    \item They are the backbone of parallel communication for sending IQ samples to a touchscreen (PortaPack) instead of going over USB.
\end{itemize}

\subsection*{P22}
\begin{itemize}
    \item Clock header and bus expansion signals.
    \item This is where the \textbf{0.5 ppm TCXO} modules are installed. When inserted, the primary clock changes to the one from this header, stabilizing the general PLL.
    \item Also exposes alternate CLKIN (pin 2), enabled with \texttt{hackrf\_clock -c p22}.
\end{itemize}

\textbf{Usage note:} Do not exceed the LPC43xx current thresholds on any pin. Exceeding them will permanently damage the central digital logic.

\end{document}
